\section{Unconditional convergence of the hybrid}
\label{sec:convergence}

The hybrid's correctness does not depend on the unresolved accelerated
locality bound.  Once AESP is stopped after any finite number of stages, local
SOR is a descent method for the original strongly convex quadratic.

\begin{lemma}[Exact objective decrease of one relaxed coordinate update]
\label{lem:coordinate-decrease}
Let \(x\in\R^n\), choose a coordinate \(u\), and let
\begin{equation}
    x^+
    =x-\frac{2\omega}{1+\alpha}\nabla_u f(x)e_u,
    \qquad 0<\omega<2.
    \label{eq:general-coordinate-update}
\end{equation}
Then
\begin{equation}
    f(x)-f(x^+)
    =\frac{\omega(2-\omega)}{1+\alpha}\bigl[\nabla_u f(x)\bigr]^2.
    \label{eq:exact-coordinate-decrease}
\end{equation}
\end{lemma}

\begin{proof}
Because the graph is simple, \(Q_{uu}=(1+\alpha)/2\).  For a scalar step
\(d\), quadratic expansion gives
\[
    f(x+de_u)=f(x)+d\nabla_u f(x)+\frac12Q_{uu}d^2.
\]
Substitute
\(d=-\omega\nabla_u f(x)/Q_{uu}
=-2\omega\nabla_u f(x)/(1+\alpha)\) and simplify.
\end{proof}

\begin{theorem}[Finite convergence of the local SOR tail]
\label{thm:tail-convergence}
Let \(x_0\in\R^n\) be any finite handoff point, possibly with a signed
gradient.  Suppose a local queue implementation has the following fairness
property: whenever a coordinate satisfies
\(
\abs{\nabla_u f(x)}\geq\alpha\epsppr\sqrt{d_u}
\), it is eventually processed unless it becomes inactive first.  Run
\cref{eq:general-coordinate-update} with a fixed \(\omega\in(0,2)\) until no
active coordinate remains.  Then the method terminates after finite
degree-weighted work and returns \(x_{\rm out}\) satisfying
\begin{equation}
    \norm{D^{-1/2}\nabla f(x_{\rm out})}_\infty<\alpha\epsppr.
    \label{eq:tail-final-certificate}
\end{equation}
Moreover,
\begin{equation}
    \Work_{\rm tail}
    \leq
    \frac{(1+\alpha)\,[f(x_0)-f(x^0)]}
         {\omega(2-\omega)\alpha^2\epsppr^2}.
    \label{eq:general-tail-work}
\end{equation}
\end{theorem}

\begin{proof}
Every processed coordinate is active at execution time, so
\begin{equation}
    \bigl[\nabla_u f(x)\bigr]^2\geq\alpha^2\epsppr^2d_u.
    \label{eq:active-square}
\end{equation}
By \cref{lem:coordinate-decrease}, its objective decrease is at least
\[
    \frac{\omega(2-\omega)\alpha^2\epsppr^2}{1+\alpha}d_u.
\]
Summing over all executed updates and using \(f(x)\geq f(x^0)\) yields
\cref{eq:general-tail-work}.  Since all degrees are positive integers, finite
work implies finitely many executed updates.  Fairness then implies that at
termination no active coordinate remains, which is
\cref{eq:tail-final-certificate}.
\end{proof}

\begin{corollary}[PPR accuracy]
\label{cor:tail-ppr-accuracy}
Under the assumptions of \cref{thm:tail-convergence}, the returned vector
\(\widehat\pi=D^{1/2}x_{\rm out}\) satisfies
\begin{equation}
    \norm{D^{-1}(\widehat\pi-\pi)}_\infty<\epsppr.
    \label{eq:tail-ppr-accuracy}
\end{equation}
\end{corollary}

\begin{proof}
Apply the source certificate \cref{eq:ppr-certificate}.
\end{proof}

\begin{theorem}[Unconditional convergence of the two-phase method]
\label{thm:hybrid-convergence}
Assume every AESP inner stage is solved to the prescribed inexactness
condition and the AESP phase is stopped after a finite stage \(J\).  Initialize
the tail from \(x^{(J)}\), reset momentum, and run the fair local SOR method of
\cref{thm:tail-convergence}.  Then the full hybrid terminates with the PPR
accuracy guarantee \cref{eq:tail-ppr-accuracy}.
\end{theorem}

\begin{proof}
The AESP phase consists of finitely many well-defined local inner solves and
therefore returns a finite vector \(x^{(J)}\).  The tail conclusion follows
from \cref{thm:tail-convergence,cor:tail-ppr-accuracy}.
\end{proof}

\begin{remark}[What convergence does not prove]
The theorem establishes correctness and finite work.  For an arbitrary
handoff, the quantitative bound in \cref{eq:general-tail-work} has
\(1/\epsppr^2\) dependence.  Acceleration requires a handoff point with a
sufficiently small objective gap or weighted residual mass; this is addressed
next.
\end{remark}
